\section{Conclusion}
\label{sec:conclusion}

This paper presented vLLM-FI, a fault-injection and reliability-evaluation framework integrated with the native vLLM runtime. The framework extends fault modeling to account for computational workload and inter-device communication, coordinates injection across distributed Workers, and performs in-place bit flips on GPU-resident tensors. Together with support for quantized inference and integration with \texttt{lm-evaluation-harness}, these capabilities enable systematic reliability evaluation across downstream tasks and parallel execution configurations.

Experiments demonstrate both the consistency and efficiency of vLLM-FI under the evaluated settings. Its fault-injection correct rates differ from MRFI by at most 0.62 percentage points, while its end-to-end runtime achieves a $1.81\times$--$2.41\times$ speedup. With the inference runtime held fixed, the CUDA backend achieves up to a $5.22\times$ speedup in total experiment runtime over an MRFI-style CPU-mediated backend under compute-fault injection. The reliability study further shows that task vulnerability depends on the fault model, longer generation increases sensitivity to persistent weight corruption, and parallel execution introduces distinct communication and stage-dependent fault effects. These findings motivate evaluation across multiple tasks and deployment configurations rather than reliance on a single benchmark. By reducing experimental cost and connecting distributed fault injection with standardized task evaluation, vLLM-FI provides a practical foundation for studying LLM reliability and informing selective protection and deployment decisions.
